\section{Introduction}
\label{sec:introduction}

Language-model agents build solutions over a sequence of interactions: they retrieve evidence, invoke tools, inspect observations, and revise their next actions. Frameworks such as ReAct make this interaction between reasoning and environmental feedback explicit \citep{yao2023react}, while memory-management systems address the difficulty of maintaining useful information over extended sessions \citep{packer2023memgpt}. The accumulated history is therefore part of the agent's working context. Keeping this context available incurs a growing cost: autoregressive Transformers store key--value (KV) representations of previous tokens to avoid recomputing them during decoding, and this cache grows with the sequence length. Efficient cache management is consequently central to serving long interactions under a fixed memory budget \citep{kwon2023pagedattention}.

The compression problem is not determined by context length alone. A tool output may contain a fact needed several actions later, and an answer may depend on evidence spread across different observations. Recency and relevance to the current question can therefore provide incomplete guidance about what to retain. Query-agnostic compression addresses this setting by constructing a compressed cache before a particular downstream question is known \citep{kim2025kvzip}. The objective is to preserve information that the language model can use across subsequent requests, while removing cache entries whose contribution does not justify their memory cost.

KVzip estimates this contribution by asking the underlying language model to reconstruct the context and measuring how strongly cached entries are attended to during reconstruction \citep{kim2025kvzip}. Fast KVzip learns a lightweight gate to predict these scores from hidden representations, avoiding reconstruction when scoring a new context \citep{kim2026fastkvzip}. Its gate operates separately on each token representation. These representations already contain context from the frozen Transformer, but the selector itself does not learn an additional exchange of information across token positions. This motivates examining whether explicit relational modeling within the selector can better support cache retention decisions.

We develop a graph-based cache selector that augments Fast KVzip with learned message passing within context blocks. Tokens form the nodes of a graph whose edge weights depend on their hidden representations. Separate graph parameters for each layer and KV head allow the selector to represent different relationships within the same context. The resulting messages enrich the representations used to predict retention scores; the retained KV vectors themselves remain unchanged. Graph-based eviction has also been explored through similarity-based score propagation \citep{li2025graphkv}. Our focus is on learning the relational scoring function and adapting its decisions to the behavior of the frozen language model.

Training has two stages with distinct roles. Reconstruction-score supervision first serves as a warm-up, providing direct targets for individual cache entries and an informed starting point for the selector. Answer supervision then adapts this selector to downstream generation: given a context and question, the full-cache model supplies a reference answer, and the selector is trained so that the same model can predict that answer from a compressed cache. The question supplies training feedback without becoming an input to the context selector. This separates a broadly available importance signal from an objective tied to how the model uses the retained information.

\medskip
\noindent\textit{[Final introduction paragraph reserved for the principal empirical findings and the paper's main contributions.]}
